\documentclass{../UTNetLabFa}

\title{کلیات شبکه و پشته‌ی \textLR{TCP-IP}}
\labnumber{۰}

\begin{document}

% \tableofcontents
% \listoffigures
% \listoftables
% \lstlistoflistings
%\listofcodes
\pagebreak

\part{تعاریف}
\section{شبکه}
\section{اینترنت}
اینترنت یک سیستم جهانی است که از اتصال میلیون ها کامپیوتر در سراسر جهان ساخته شده است.

برای تنظیم کردن مسیریاب از دستورات زیر استفاده کنید:

\begin{figure}[h]
	\centering
	{\selectlanguage{english}\includegraphics[height=64pt]{../img/fanni}}
	\caption{\FR{تصویر دانشکده‌ی فنی - دانشگاه تهران}}
	\label{fig:fanni}
\end{figure}

\subsection{فونت‌ها}
نمونه فونت‌های قابل استفاده:
\begin{table}[h]
	\label{tab:farsi-fonts}
	\caption{جدول فونت های موجود}
	\centering
	\begin{tabular}{|l|l||R|R|}
		\hline
		\bfseries \LR{Font name} & \bfseries\en{Command}           & \bfseries نام & \bfseries    نمونه متن            \\ \hline\hline
		\en{Nazli}               & \en{\textbackslash{}textnazli}  & نازلی         & \textnazli{یک متن نمونه از فونت}  \\ \hline
		\en{Nazli Bold}          & \en{\textbackslash{}textnazlib} & نازلی درشت    & \textnazlib{یک متن نمونه از فونت} \\ \hline
		\en{Titr Bold}           & \en{\textbackslash{}texttitr}   & تیتر          & \texttitr{یک متن نمونه از فونت}   \\ \hline
		\en{Homa}                & \en{\textbackslash{}texthoma}   & هما           & \texthoma{یک متن نمونه از فونت}   \\ \hline
	\end{tabular}
\end{table}



\begin{otherlanguage}{english}
	\section{English Section}
	This is a english section.%\footnote{\en{sample footnote}}
	%	\subsection{English subsection}
	Subsection text.
	%	\subsubsection{English SubSubSection}
	Subsubsection text.
\end{otherlanguage}

\part{ادوات شبکه}
\section{اتصالات}
در این بخش توضیحاتی در رابطه با ادوات شبکه قرار می‌گیرد.
\subsection{سیمی}
انواع ارتباطات سیمی

\subsection{بی سیم}
اتصالاتی همانند \textLR{WiFi} و \textLR{Bluetooth}.
\subsection{نوری}
اتصالات فیبر نوری به علت نویز کمتر در مسیر، امکان ارسال به مسیرهای دور را با سرعت بالا و خطای کم را فراهم می‌کنند.
\begin{question}
	\item متن سوال
\end{question}

\begin{report}
	\item نمونه گزارش
\end{report}

\section{دستگاه‌ها}

\subsection{هاب}
\subsection{سوئیچ}
\subsection{مسیریاب}
\subsubsection{تنظیمات مسیریاب}
{
	\selectlanguage{english}
	\begin{lstlisting}[language=cisco,caption=\FR{تنظیمات مسیریاب}]
config term
	no ip routing
	bridge 1 protocol ieee ! for STP protocol
	int f0/0
		ip addr 128.238.61.1 255.255.255.0
		bridge-group 1
		no shut
		exit
	int f0/1
		ip addr 128.238.61.2 255.255.255.0
		bridge-group 1
		no shut
		end
	Ctrl+Z
\end{lstlisting}
}
\subsection{سوئیچ های نسل جدید}
\subsubsection{کنترل کننده}
\subsection{پایشگر (\en{FireWall})}

\part{شبیه سازی}


\end{document}
